% ============================================================
% Part III — Desktop Section — M0 (Introduction) — DRAFT v2 (2026-07-20)
% RESTRUCTURED to the author's six-move introduction structure (FCDP §10.1;
% memory feedback-user-intro-structure): hook > context > disruption >
% resolution > signpost > thesis. v1 preserved as DESKTOP-M0-DRAFT-v1.tex
% (major iteration on a copy).
% Move map: hook = NEW (headset shelf + two K&S verbatims) · context = M0.1 +
%   NEW frame passage + M0.2 · disruption = NEW (rhetorical question) ·
%   resolution = NEW (gap + frame-as-resolution, methodology folded) ·
%   signpost = M0.4 (GOLD STANDARD roadmap, untouched) · thesis = M0.3
%   (opener retrofit only).
% Conventions carried from v1 (MLA; ("Phenomenological Evaluation" NN);
%   first-person for K&S work; "I" for dissertation-argument claims; bans:
%   "record", sentence-initial "And", metatextual signposting, gambling
%   vocabulary, bare "incorporation" for our theory [always "rhetorical
%   incorporation"; bare form = Calleja's theory only]).
% Verified quote pages (PDF): indisputable-presence 390; not-a-single-student
%   389; "a hassle" (results wording) 388; hypothesis + unlikely 383;
%   restraint 390; no-collaboration 389; figures 388; reasons 389; scarcity
%   381; phenomenological framing 387; Heim AWS 182/185; Radianti 1/22.
% NO-STALL SWEEP APPLIED 2026-07-29 (user-ordered at assembly; doctrine
%   postdated M0 approval): "stall at any station" -> "fall silent at any
%   station" (frame passage); "lets it stall" -> "falls silent before it"
%   (resolution); thesis "a stall at the passage" -> "the world falls
%   silent at the passage" (M7.2's approved language); "completion
%   stalled" -> "the act was withheld" (thesis close).
% DOCTRINE (user 2026-07-31, readthrough batch 1): never phrase the
%   A3->A4 diagnosis as "the space provides no act" or "students get
%   stalled at A3" — a potentiality to actualize is always present
%   (watching is itself an actualization); the failure is in the WORTH of
%   the potentialities the world furnishes. Design's job = make them
%   worthwhile.
% REVISION-PASS MARKERS: (a) social-interior/Second Life passage — revise
%   against the new-media/presence sources once ingested + indexed (agent
%   running); (b) roadmap part-count + thesis — re-verify at first complete
%   section draft; (c) USER NS 2026-07-31: once the section's argument and
%   content are finalized, revise the introduction (esp. M0.3 thesis) and the
%   close (M7) so the thesis is clear and definitively supported.
% ============================================================

\section{From the Pilot to the Deployment}

% === (I) HOOK — [REVIEW: v3 this round; K&S verbatims added per user note.
%     Full-citation/convention footnote RELOCATED here from M0.1.]

In the spring of 2022, the science library at the University of California,
Irvine kept a shelf of Meta Quest 2 headsets on loan to the students of a
biomedical-engineering elective.\footnote{Part III draws on several of my
previous publications; the article referenced here---which is reproduced in
Appendix~\ref{app:vr-pilot}---is the following, and will be cited
parenthetically by abbreviated title and page number hereafter: King, Christine
E., and Dalton Salvo. ``Phenomenological Evaluation of an Undergraduate Clinical
Needs Finding Skills Through a Virtual Reality Clinical Immersion Platform.''
\textit{Biomedical Engineering Education}, vol.~4, 2024, pp.~381--97,
https://doi.org/10.1007/s43683-024-00139-5.} Any student could borrow one at no
cost and, through it, stand inside an operating room---a space few
undergraduates will ever be admitted to. It is, we would later conclude,
``indisputable that a majority of our students felt that VR increased their
sense of presence and immersion'' (``Phenomenological Evaluation'' 390). Yet
``[w]hile we had several Meta Quest 2's housed in the school's library, not a
single student went there to check one out'' (389). That untouched shelf is this
section's problem in miniature: a virtual world can offer presence and still
lose its students to the more convenient copy of itself. What follows examines
that problem at scale---three quarters of the same course, redeployed on
ordinary computers to 241 students across the University of California---and
asks what a virtual learning environment (VLE) must offer, beyond presence, to
be worth entering.
% [Hook follow-up APPROVED 2026-07-20.]

% === (II) CONTEXT — part 1: the pilot (M0.1, locked text; seam edits marked)

Clinical immersion for non-medical staff within operating rooms and other
restricted medical spaces is scarce by structure. The programs that admit
undergraduates to such clinical settings draw substantially more applicants than
they can serve, so that a majority of the students who seek a clinical-immersion
experience will not have one---an inequity the biomedical-engineering literature
documents against a clear demand.\footnote{Guilford, Kotche, and Schmedlen,
surveying clinical-immersion experiences across biomedical-engineering programs,
find that only 48\% of applicants to clinical-immersion courses and 24\% of
applicants to immersion programs are able to participate annually (113--22).
Moravec et al.\ report on one such program---a short-term clinical immersion for
undergraduate biomedical engineering students---and the best practices such
programs have developed under the same constraint of scarce placements and clear
demand (217--23).} The virtual clinical-immersion platform whose deployment this
section reads was built against that scarcity: filmed operations, physician and
engineer interviews, and explorable procedure rooms, delivered so that any
enrolled student could be emplaced where only a limited few traditionally have
access. In its first published test, a pilot cohort of twenty-two
% [SEAM EDIT: "Its first published test came in the spring of 2022, when a
%  pilot cohort..." — the hook now carries spring 2022; wording adjusted.]
met the same operating-room footage in four competing formats: 2D video on
YouTube, an interactive computer build,\footnote{The interactive computer build
was developed in Unreal Engine~4 for standard PCs: essentially the same
explorable virtual environment as the Meta Quest 2 build, but navigated by mouse
and keyboard on a flat monitor, and delivering the traditional 2D videos rather
than the stereoscopic footage (``Phenomenological Evaluation'' 383).} a Google
Cardboard head mounted display (HMD),\footnote{Google Cardboard is a minimal
head-mounted display: a folded cardboard or plastic viewer holding two lenses,
into which the user inserts a smartphone. The phone's screen, split into a
stereoscopic pair and viewed through the lenses, produces a three-dimensional
image, and the phone's own motion sensors track the rotation of the head. In
this course the Cardboard route delivered the operating-room footage through
YouTube's VR mode---video only, with none of the interactive environment of the
PC and Meta Quest 2 builds (``Phenomenological Evaluation'' 383). YouTube has
since discontinued this VR viewing mode, and the route the Cardboard format
relied on is no longer available. Given the nature of this setup, of all VR methods it is the most prone to producing
Alternate World Syndrome (AWS): motion sickness, physical disorientation, and
physical discomfort. In Michael Heim's definition, ``AWS is the relativity
sickness that comes from switching back-and-forth between the primary and
virtual worlds'' (\textit{Virtual Realism} 182); the VR experience ``can mix
images and expectations from alternate worlds so as to distort our perception of
the primary world, making us prone to errors in mismatched contexts'' (185).}
and a Meta Quest 2 virtual reality (VR) HMD. Interestingly, across twenty
post-course responses we discovered a two-part verdict that would not line up.
The highest level of presence\footnote{We do not stipulate a formal definition
of presence in the article; it uses the term in the sense standard to the
VR-pedagogy literature it draws on---the felt sense of being physically present
in the virtual rather than the actual environment---and treats it, alongside
embodiment, as one of the ``intensely subjective experiential phenomena'' a
phenomenological method is fitted to assess (``Phenomenological Evaluation''
382, 387). The construct is anchored to Makransky and Lilleholt (2018) and
Makransky, Terkildsen, and Mayer (2019) and measured throughout by
self-report.} went to the headset: 70\% reported that the Meta Quest
VR version elicited a greater felt sense of presence. Student preference, on the
other hand, went the other way: 40\% judged the traditional 2D YouTube
videos most educationally beneficial. The two VR HMD formats together drew only
that same share---20\% choosing the Google Cardboard and 20\% the Meta Quest
2---while 10\% selected the PC build and 20\% judged all versions equally
beneficial (``Phenomenological Evaluation''
388). The students' explanation as to why the 2D videos were perceived as being
more educationally beneficial was largely grounded in physical
discomfort\footnote{The discomfort belonged overwhelmingly to the Google
Cardboard rather than the Meta Quest 2: asked which version was \textit{least}
educationally beneficial, 40\% of students named the Cardboard against only 5\%
for the Quest (``Phenomenological Evaluation'' 388). Among the students who
judged the 2D videos most beneficial, 38.9\% cited physical
discomfort---disorientation, the weight of the headset over long viewing
periods, the difficulty of wearing glasses inside an HMD---a figure that pools
the two VR formats (389). Our own judgment in the article is that some of these
issues ``would have been less relevant'' had students used the Meta Quest 2
rather than the Cardboard, ``which is a relatively horrible experience by
comparison'' (390).} and convenience.\footnote{Among the students who judged the
2D videos most beneficial, 44.5\% named ease of use as the ground: convenience
and accessibility of access, and the quality-of-life controls the VR formats
lacked---variable playback speed, better volume and playback control, and the
ability to zoom in on the footage (``Phenomenological Evaluation'' 389). The
rest of the students' reasoning is broken down as follows: 11.1\% cited the
inability to take notes while wearing an HMD, and 5.6\% had never tried a VR
version at all (389).} The loaner headsets went unused because, as one student
described it, having to go to the library to check out the VR headset is ``a
hassle'' (388).
% [SEAM EDIT: "Despite making several Meta Quest HMDs available for checkout
%  through UCI's library, not a single student went to the library because..."
%  — the not-a-single-student fact now lives in the hook (389, verbatim);
%  this sentence trimmed to the 388 results-wording quote to avoid repetition.]

In our previous article, we stated our expectations at the start. We did not
expect virtual clinical immersion to rival physical presence in the operating
room: ``[i]t is unlikely that VR clinical immersion programs will be as
pedagogically effective as in-person observation'' (``Phenomenological
Evaluation'' 383). We hypothesized instead that the first-person perspective
enabled by VR, ``and the enhanced sense of presence and agency it generates, has
the potential to be more beneficial to students' development of such [clinical
needs-finding] skills than 2D videos as it more closely simulates the experience
of being physically present in the operating room'' (383). The results
substantiated half of that initial hypothesis---presence was at its highest in
the Meta Quest VR environment, but student preference did not follow suit---and
on the perceived pedagogical benefit, we were unable to render a verdict.
Because the course's commitment to open access let every student default to the
most convenient format,\footnote{Other than the occasional in-class use of the
VR headsets, students were not required to use any particular format and were
allowed to use whichever of the four versions of the video content they wanted:
YouTube, the PC version, Google Cardboard, or Meta Quest 2.} the four versions
could not compete on equal grounds: while a majority of students undeniably felt
more present in virtual reality, we could not ``conclusively claim that
interactive VR learning environments are more or less pedagogically efficacious
compared to 2D videos'' (390).\footnote{Our inability to render a verdict places
the pilot within a larger evidential pattern. Comparative studies that hold
content constant repeatedly find that added immersion raises felt presence
without raising learning: Makransky, Terkildsen, and Mayer, porting an identical
science simulation from desktop to headset, measured more presence and less
tested learning; Parong and Mayer's content-matched comparison found a slideshow
outperforming immersive VR on factual learning even as VR won every affective
measure; and Mayer, Makransky, and Parong's review of two decades of media
comparisons finds no reliable learning advantage for immersive VR as a medium.}
The pilot thus established two aspects and clearly drew a separation between
them: presence was achieved to a much higher extent in VR, but was not enough to
override the students' preference for convenience when having to view the
educational content at home on their own time. Whether the interactive virtual
world was more pedagogically effective---whether felt presence converts into
pedagogical worth---was left standing as a question our pilot's own conditions
could not answer.

% === (II) CONTEXT — part 2: the frame (NEW passage, approved 2026-07-20 with
%     "rhetorical incorporation" terminology fix)

Three concepts, built across the preceding parts, do the analytical work here.
The first is the actualization chain: engagement understood as the actualization
of desire, running from perception (A\textsubscript{1}) through the charged
image (\textit{phantasma}, A\textsubscript{2}) and committed judgment
(\textit{doxa}, A\textsubscript{3}) into completed action
(A\textsubscript{4})---a sequence a designed world can solicit, sustain, or
fall silent at any station. The chain also has a zero point
(A\textsubscript{0}): before any of its stations can fire, the student and the
world must reach one another at all---the threshold condition every later
moment presupposes. The second is rhetorical incorporation: the union a virtual world
achieves with its inhabitant, one in substance but dual in
being---world-disclosedness, the world opened to a student as a place they are
in, and co-disclosedness, the same world opened to students as a place they are
in together. The third is veri(dis)similitude: the requirement that a virtual
world be similar enough to what it renders to be entered as real, yet different
enough from its source material to justify its existence. A virtual learning
environment, in these terms, is not a delivery medium but a manufactured and
curated rhetorical ecology---an environment authored, in every element, toward a
single pedagogical end, whose design succeeds when both registers of disclosure
open and rhetorical incorporation is achieved.

% === (II) CONTEXT — part 3: the institutional resolution (M0.2, locked text)

After the pilot, we moved the course to a 2D-PC primary
deployment; while students cannot experience stereoscopic 3D video without a VR
headset, the videos in this interactive space are still 180- and 360-degree
presentations. Our primary reason for moving away from the 3D stereoscopic VR
development was due to logistics, feasibility, and necessary scale. This class
needed to be made available to undergraduate students across the entire
University of California (UC) system, with class sizes up to 120 students each
offering. Acquiring and distributing 120 Meta Quest headsets to each student
every quarter was simply not possible. Secondly, we chose to prioritize the
Windows PC-based version first as it was the largest student platform population
(compared to Apple laptops).\footnote{While we prioritized development for the
PC version, we are currently in the process of developing a version for Apple
computers, and development for VR headsets will commence once the Apple version
is operational.}

As just noted, this class (******) is offered systemwide to undergraduate
biomedical engineering (BME) students across the UCs.\footnote{This class is
available to undergraduate BME students at every UC except the University of
California, San Francisco (UCSF), as UCSF is exclusively graduate-level health
sciences.} Students arrive with different course schedules and even different
academic calendars (Berkeley and Merced on semesters while the rest of the
universities are on quarters), so a fully asynchronous online format is not a
convenience of choice, but the condition of the class's systemwide existence, as
a synchronous version would be nearly impossible to logistically organize. An
in-person version was never possible at all, and cohort-wide scheduled
simultaneity was never in the design space: the deployment maximized the only
availability there was.
% AUTHOR FOLLOW-UP (******): (1) class name/number + catalog details from the
% instructor; (2) how is enrollment handled for Berkeley/Merced students given
% the semester calendar?

These constraints are more than administrative history; they define what kind of
thing the virtual learning environment is. The virtual world was built, in the
first place, because the actual clinical world admits so few---the operating
room's structural scarcity was the problem, and universal access was the
solution the virtual environment promised. A virtual world that students cannot
enter would therefore fail in a specific and instructive way: it would
reproduce, in software, the very exclusion it was built to overcome.
Availability is thus not a logistical afterthought but the primary pedagogical
property of this course and VLE: before a virtual world can show, teach, or
immerse, it must be enterable, and whatever a student cannot enter can disclose
nothing to them. In the vocabulary of the actualization chain, this is a
condition at A\textsubscript{0}, the chain's threshold: perception, the charged
image, committed judgment, and the act are all stations a student can only
reach inside a world they occupy, so a world that cannot be entered does not
engage its students poorly---it never begins to engage them at all.
Availability is the threshold's deployment face: the form A\textsubscript{0}
takes when the world must be downloaded, installed, and run before it exists
for its student. Nor is the scaling wall we hit peculiar to this course.
Immersive VR in higher education remains largely at the prototype
stage---experimental and development work rather than a regular part of teaching
(Radianti et al.\ 1, 22)---a pattern that will be discussed in greater detail
later. One further consequence of the deployment should be noted now and
developed later: a fully asynchronous format changes the time of the course.
There is no shared class hour; each student meets the virtual world on their own
schedule and at their own pace. The VLE itself is built for company---it is
multiplayer, with spatial proximity voice chat---and we clearly inform students
of this capability and encourage them to meet there in their assigned groups.
But with no scheduled hour to gather them, convening is left to each group's own
initiative, and the default is that a student enters the virtual hospital alone.
Whether students take up the invitation to meet---and what solitary viewing does
to their experience of the virtual world when they do not---are questions the
three terms of survey responses will answer.

% === (III) DISRUPTION (NEW; approved 2026-07-20 with largest-share
%     parenthetical added)

Set side by side, the pilot and the field leave a precise question standing.
70\% of our pilot cohort felt more present in the headset, yet 40\% judged
plain 2D video the most educationally beneficial format (the
largest share of the cohort to settle on any single format), and the loaner
headsets sat untouched in the library; the comparison literature, holding
content constant across media, finds the same shape at scale---added immersion
reliably raises presence and just as reliably fails to raise learning, while
presence itself proves attainable on an ordinary monitor. Presence, in short, is
neither scarce nor sufficient. What, then, must a virtual learning environment
add to presence before the world it renders earns its place in a curriculum?
What would a student gain inside the virtual world that the same footage, one
click away on YouTube, cannot give them---and what kind of analysis could even
name that difference?

% === (IV) RESOLUTION (NEW; approved 2026-07-20 with Second Life rephrase)
% REVISION-PASS MARKER: the social-interior passage (from "There is a material
% reason" through "authored for a curriculum") is to be revised against the
% new-media/presence sources once ingested + indexed (Boellstorff, Davis &
% Boellstorff, Champion, Spagnolli, Mantovani & Riva, Riva, et al.).

The VR-pedagogy tradition has measured presence with increasing precision while
conceding it cannot say what presence is for; the social-presence tradition has
known since Short, Williams, and Christie's founding work in 1976 that the felt
company of others carries satisfaction and learning, but it matured in distance
education---discussion forums, video sessions, text---around media that were
channels rather than places. The two research programs descend from the same
1976 origin and cite one another almost never, having built their houses in
different fields around different objects. There is a material reason the
divide has persisted: until recently, few deployed pedagogical virtual
environments had a social interior. Multiplayer capability and spatial
proximity voice chat---the features that make a virtual world a place where
students can be together---have been rarely implemented. Where educators
have taught inside socially capable virtual worlds in relatively recent
times---the Second Life classrooms documented by anthropologists and education
researchers offer one example---the world has often been a general-purpose
social platform adapted to teaching rather than an environment authored for a
curriculum.\footnote{On the education-research side, the fullest example in
this section's literature review is Merchant et al.'s structural-equation study
of college chemistry instruction inside Second Life (N\,=\,204), taken up in
detail below. On the anthropological side, Tom Boellstorff---an anthropologist
at UC Irvine---conducted multi-year ethnographic fieldwork inside Second Life
through the avatar Tom ``Bukowski,'' the fieldwork behind his monograph
\textit{Coming of Age in Second Life}; he recounts that work in ``Virtual
Worlds and Futures of Anthropology'' and carries the program forward to the
present platform generation in ``Toward Anthropologies of the Metaverse.''} A
desktop VLE of the present kind therefore sits at a junction
neither literature has occupied: an explorable virtual hospital authored for one
clinical curriculum, deployed for credit across a university system, its
students group-assigned and equipped with in-world proximity voice chat---at
once a
rendered world of the kind the VR-pedagogy literature measures and a social
space of the kind the social-presence literature theorizes. The VR-pedagogy
studies compare media and score presence but do not ask what being-together
inside the world does; the social-presence studies measure felt community in
forums and video sessions but not inside a navigable world; to our knowledge, no
analysis has read a single deployment with both sets of questions in hand.
Reading one requires a frame built for the junction, and that is what this
dissertation's rhetorical-ontological and phenomenological frame supplies. Where
the field's instruments report that presence rose or interest fell, they cannot
say how a wholly authored environment solicits a student's desire, carries it
toward action, or falls silent before it; the actualization chain names those
moments.
Where presence names a feeling, world-disclosedness and co-disclosedness name
the two registers in which a virtual world is opened to its inhabitants---as a
place I am in, and as a place we are in together---so that what a design
achieves or forfeits can be located, not merely scored. A VLE, on this reading,
is a manufactured and curated rhetorical ecology with a single pedagogical end,
and its effectiveness is a design question in the strictest sense: the
environment teaches when it is designed so that both registers open---when
rhetorical incorporation is achieved.

% === (V) SIGNPOST (M0.4, GOLD STANDARD — untouched; see FCDP §10.2)

The survey given to all students across the Winter 2025, Spring 2025, and
Winter 2026 quarters was designed, fielded, and collected by myself and
Prof.\ Christine King---the instructor of record---and is read here
through the rhetorical-ontological and phenomenological frame this dissertation
has been building. The data yield the findings, with every count and quoted
response coming from the survey corpus. The theoretical frame supplies the how
and the why: the actualization chain and the methodology of
Rhetorical-Ontological Diachronic Analysis (RODA) serving as a background
grammar for interpreting students' responses---which moment of engagement is in play when a
student reports presence, confusion, the flight to YouTube, or the desire for
more socially interactive course design. Accordingly, what follows moves through
seven parts. The first reviews the relevant literature in its two traditions:
the desktop route as a pedagogical modality in its own right, with presence
relocated from hardware to the design of the environment; the comparative
evidence that more presence does not deliver more learning; the rationale for
asynchronous deployment at scale; and the research on learning together in
shared virtual space---social presence and computer-supported collaborative
learning---whose near-total lack of contact with the VR-pedagogy literature is
itself a finding this section works across. The second part describes the VLE
and its deployment: the virtual hospital as a designed world, the cohort of 241
students across the three quarters, and the survey instrument together with the
standing limits of what it can show. The third presents the survey findings in
order: the threshold problems that kept some students out of the virtual world;
what students found, and could not do, once inside; the flight to YouTube; the
presence students attested on ordinary screens; the group co-watching capability
that was had, wanted, and largely unactivated; and the shift in student
skepticism from technical friction to the question of the platform's value. The
fourth part conducts the rhetorical-ontological and phenomenological analysis:
it locates the primary rupture point in the student experience, takes up the
justification question the data raise, and shows how the VLE is world-disclosed
and co-disclosed to its students. The fifth turns diagnosis into design, ranking
the design levers by the moment of engagement each one works on, with the
activation of the WE given the fullest treatment. The sixth states the limits of
what this survey data can unconceal. The seventh names what the analysis hands
forward---to the boredom analysis, a separate section that follows this one, and
to the field.
% NOTE (user 2026-07-20): the boredom analysis is its OWN SECTION, not a
% subsection of this one. RE-VERIFY the seven-part count at first complete
% section draft.

% === (VI) THESIS (M0.3 relocated to final position; opener retrofit marked)

By reading the three quarters of survey data through this frame, the following
analysis will substantiate a claim with three connected parts.
% [SEAM EDIT: was "The claim this section will substantiate can be stated at
%  the outset, in three connected parts." — "at the outset" no longer true in
%  final position; retrofit per approved plan.]
First: the deployed desktop world succeeds where our pilot said screens could
succeed. Across three terms, students attest to presence on ordinary monitors at
rates matching what the pilot measured in VR headsets; a virtual world on a flat
screen can still place its student somewhere, and world-disclosedness does not
require a headset. But presence, I will show, is not the finish line. As
deployed, the virtual hospital gives students a great deal to see---filmed
operations, physician and engineer interviews, explorable procedure
rooms---and watching is nearly all it gives them to do. Watching is itself an
act: a student can stand in the virtual operating room and refuse to attend,
and the one who watches has actualized a potentiality the world furnished. The
shortfall lies in what else the world makes available: no task to carry out
inside it, no procedure to attempt, no problem whose solution it asks of the
student---no richer or more pedagogically effective potentiality for the
student's engagement to actualize. In the chain's terms, the deployment
carries engagement to committed judgment (A\textsubscript{3})---the world wins
the student's attention and their assent that it is worth taking
seriously---but the only completed act (A\textsubscript{4}) it affords that
commitment is one the student could have performed anywhere: watching. The
world falls silent at the passage from A\textsubscript{3} to
A\textsubscript{4} not because no act follows, but because the acts it affords
stop at watching. Second:
because every video housed in the virtual hospital also exists one click away on
YouTube, the desktop world carries a standing burden of justification. It must
be different enough from the video of itself to justify existing, and that
difference can only be made of what video cannot render: a place to be in, a
task to do, and others to be there with. Third: of those three materials, the
last is the one this deployment alone possesses. Watching together in the same
virtual space, within earshot of one another, is a capability no other route in
this program's history has offered---the published VR version had no real-time
collaborative features at all, and students could only enter it alone
(``Phenomenological Evaluation'' 389). Where place and task deepen
world-disclosedness, being there together opens a different register
altogether---co-disclosedness: the virtual world disclosed not to each student
alone but to a WE. I will argue that this capability is at once the desktop
world's deepest answer to ``why not just YouTube?'' and, as the data
demonstrates, its deepest unrealized asset. The scope of this claim is a matter
of development history rather than principle: nothing in a headset forbids
company; it is on the desktop route that company was built. What the analysis
owes, then, is evidence for each part: that the world was disclosed, that
the act was withheld, and that co-disclosedness---the one capability that could
justify the world---went largely unused.
% [SEAM EDIT: "What the section owes" -> "What the analysis owes" (self-
%  reference trim). NOTE: thesis to be revised once the full section exists,
%  per user 2026-07-20.]

% ============================================================
% v2 LOCKED 2026-07-20 (hook follow-up approved; seams 1-3 approved; all six
% moves user-ruled). OUTSTANDING REVISION PASSES: (a) social-interior passage
% in the resolution — revise against new-media/presence sources once indexed
% (ingestion Phase 1 done, snapshot doc-text-20260720T145233Z; entry phases
% user-gated); (b) roadmap part-count + thesis — at first complete section
% draft; (c) class-name (******) + Berkeley/Merced enrollment follow-up.
% ============================================================
